\documentclass[11pt,letterpaper]{article}
\usepackage[margin=0.88in,headheight=15pt]{geometry}
\usepackage[T1]{fontenc}
\usepackage{lmodern,microtype,amsmath,amssymb,amsthm,mathtools}
\usepackage{booktabs,tabularx,array,xcolor,fancyhdr}
\usepackage[colorlinks=true,linkcolor=blue!45!black,urlcolor=blue!50!black,citecolor=blue!45!black]{hyperref}
\definecolor{ink}{RGB}{28,49,68}
\pagestyle{fancy}\fancyhf{}
\setlength{\parindent}{0pt}\setlength{\parskip}{5pt}
\setlength{\emergencystretch}{2em}
\setcounter{tocdepth}{1}\numberwithin{equation}{section}
\theoremstyle{plain}
\newtheorem{theorem}{Theorem}[section]
\newtheorem{lemma}[theorem]{Lemma}
\newtheorem{proposition}[theorem]{Proposition}
\newtheorem{corollary}[theorem]{Corollary}
\theoremstyle{remark}\newtheorem{remark}[theorem]{Scope}
\DeclareMathOperator{\Var}{Var}\DeclareMathOperator{\Cov}{Cov}
\DeclareMathOperator{\Tr}{Tr}\DeclareMathOperator{\Vol}{Vol}
\DeclareMathOperator{\supp}{supp}\DeclareMathOperator{\diver}{div}
\DeclareMathOperator{\Ad}{Ad}
\newcommand{\E}{\mathbb E}\newcommand{\R}{\mathbb R}
\newcommand{\dd}{\,\mathrm d}\newcommand{\OH}{\mathcal O_H}
\newcommand{\norm}[1]{\left\lVert#1\right\rVert}
\newcommand{\abs}[1]{\left|#1\right|}
\newcommand{\ind}{\mathbf1}\newcommand{\cC}{\mathcal C}
\newcommand{\cT}{\mathcal T}\newcommand{\eps}{\varepsilon}
\title{Connected estimates for selected operations}\author{}\date{}
\providecommand{\mathscr}{\mathcal}
\begin{document}\maketitle
\section{The unchanged source and the exact scalar bookkeeping}
\label{sec:source}

Let $G=SU(N)$ with normalized Haar measure, fine spacing $a$, and
\begin{equation}
s_p=N-\Re\Tr U_p,\qquad S_0=\sum_p s_p,\qquad
\rho_W=e^{-\beta S_0},\qquad Z_\beta=\int\rho_W\dd\lambda.
\end{equation}
Use the original invariant directions
$X_{e,A}F=\left.\frac{\dd}{\dd t}F(\ldots,e^{tT_A}U_e,\ldots)\right|_{t=0}$.
Write $F_i=X_iS_0=\sum_{\mu\ne\nu_i}F_{i;\mu}$, and retain every
original test shift and hypercubic average in
\begin{align}
D[f]&=\sum_i f_i\sum_{\mu\ne\nu_i}F_{i;\mu}^2,
&P[f]&=\sum_i f_iF_i^2,
&\ell[f]&=\sum_i f_iX_i^2S_0,\notag\\
\OH(f)&=\sigma\{D[f]-P[f]/3\},
&\sigma&=\gamma^{-2}a^{-2}.
\label{eq:H}
\end{align}
Here $\gamma$ is the saved source coupling. It equals the birth coupling when
the primitive is first inserted; it is not reset to a later native coupling.
For the action coefficient $\beta_s$ on the source lattice, write
$\rho_s=e^{-\beta_sS_0}$; at the original bare birth $\beta_s=\beta$ and
$\rho_s=\rho_W$. The original decomposition is
\begin{align}
V_H(f)&=-\frac{\sigma}{3\beta_s}\sum_i f_iF_iX_i,
&t_H&=-\frac{\sigma}{3}\{P[f]-\beta_s^{-1}\ell[f]\},\notag\\
r_H&=\sigma D[f]-\frac{\sigma}{3\beta_s}\ell[f],
&\OH&=t_H+r_H.
\label{eq:Hdecomp}
\end{align}
For $D_V\nu=-\diver(\nu V)$, the bare smooth identity is
$t_H\rho_s=D_{V_H}\rho_s$.
The Casimir identity gives the three ordinary rows
\begin{equation}
r_H=\sigma D[f]
-\frac{\sigma C_FN}{3\beta_s}\sum_p f_p
+\frac{\sigma C_F}{3\beta_s}\sum_p f_ps_p.
\label{eq:Casimir}
\end{equation}
The last term is the field-dependent action row; the middle term is an absolute
scalar. Both remain. The incidence sum defining $f_p$ is the original one.

For commuting square-zero source labels $\eta_i$, the general inherited source
action is
\begin{equation}
A(\eta)=\sum_{\varnothing\ne I}\eta_I A_I,
\qquad \eta_I=\prod_{i\in I}\eta_i.
\label{eq:finiteaction}
\end{equation}
All exponentials and logarithms are finite algebraic operations. For two labels
write $A_1=O_1$, $A_2=O_2$ and keep $A_{12}$. If an old exact decomposition
$O_i=t_i+r_i$ is used, its complete compensated coefficient is
\begin{equation}
b_{12}=A_{12}+\tfrac12(V_1t_2+V_2t_1)+V_1r_2+V_2r_1.
\label{eq:b12}
\end{equation}
Equation~\eqref{eq:b12} includes the responses of the action and scalar rows
in~\eqref{eq:Casimir}. The later estimates apply to the full packet, not just
one convenient piece of this expression.

Let $T_z$ denote a genuine complete component operator. If an independent
outside factor is $a(z)=e^{s(z)}/D_{\rm orig}(z)$, the actual raw coefficients
are
\begin{align}
N_0&=aZ_z,&N_i&=aT_zO_i,&
N_{12}&=aT_z(A_{12}+O_1O_2),\notag\\
N_{12}-\frac{N_1N_2}{N_0}
&=a\{T_zA_{12}+\cC_z(O_1,O_2)\},&
\cC_z(F,G)&=T_z(FG)-\frac{T_zF\,T_zG}{Z_z}.
\label{eq:rawsource}
\end{align}
Thus the independent scalar appears once in the raw connected coefficient.
Its derivative is $a'=a(s'-D_{\rm orig}'/D_{\rm orig})$ wherever that
logarithmic notation is legitimate. A factor depending on an integrated
variable stays inside the integral. Before the genuine nonzero complete body
is available, retain $N_0N_{12}-N_1N_2$ and the joint numerator without dividing
by a selected island mass.

\section{Actual admission: the part that is established}
\label{sec:admission}

The primary input is Balaban's R-II~\cite{balabanRII}, especially (1.65) and
(1.71)--(1.75). R-I~\cite{balabanRI} means the 1989 large-field basic-step
paper, reference [IV] in R-II; it is not the earlier small-field paper [I].
R-I's full defining ranges are not used here. No range margin is inferred from bibliographic metadata.

R-II's (1.65) gives native old-domain admission after its three determining-set
joins. Its (1.71)--(1.75) gives a positive real complete body and a small complex
density comparison. The body includes the internal history sum and the inherited
$V(X^-)$. Its outer pair dependence contains the quadratic $B$ form and
$V(X^-)$. These are inputs about that original operator, not a spectral-gap
assertion or a claim about a smaller marked physical core.

\subsection{A native marked bound can be carried, with its actual cost}

The original MT25a constructs a full birth-core bound of the form
\begin{equation}
b_j^{\rm birth}(\xi_j)\le C_d|\Gamma_j|A_j
\left(1+\frac{\alpha_j}{g_j}\right)^d(1+E_j(\xi_j))^{d/2},
\qquad d=2,3.
\label{eq:nativebirth}
\end{equation}
Here $\Gamma_j$ retains the actual absolute birth normalization. MT25a,
the discussion in \S100.7, and \S107.10 distinguish this native bound from
the canonical and strict-age bounds. $E_j$ is the literal birth energy; it
is not identified with a later component energy by notation. In the verified finite carrying
marked-birth class, a completed regular R endpoint has the original ordinary
pair domain. Still-live boundary coefficients have not reached that endpoint.

\begin{proposition}[Basic touched evaluation in the certified native class]
Let $F_j$ be a coefficient already certified on the original native domain
used in R-II (1.65), with its certified scalar full-core norm $b_j$. On the original
prepared range the composition $F_j\circ\Gamma^R$ is defined. On a component
with $\norm{\sigma_z}_\infty\le\eps$,
\begin{equation}
|T_z(F_j\circ\Gamma^R)|\le e^\eps Z_0b_j,\qquad
|K_z(F_j\circ\Gamma^R)|\le e^{3\eps}b_j
\label{eq:nativecarry}
\end{equation}
when the original reciprocal-body estimate applies. An insertion-level bound
$|F_j(\Gamma^R(z,\xi))|\le\widetilde b_j(\xi)$ may instead give
$e^\eps\mu_0(\widetilde b_j)$, but only after its literal pullback and energy
envelope have been checked on this actual common domain.
\end{proposition}
\begin{proof}
The native inclusion makes this particular composition legitimate. Apply the
complete-density comparison to its actual absolute bound. The raw estimate
retains the complete mass; the normalized estimate uses the complete body's
own nonvanishing bound. No source label changes the domain of a coefficient
already certified there.
\end{proof}

This is an actual, but weaker, touched admission row. It does not identify
$F_j\circ\Gamma^R$ with the stored stable physical representative
$R_j(P_{\rm old}q_{\rm old}\Gamma^R)$, and it does not replace
$(1+\alpha_j/g_j)^d$ by a saved age weight. An arbitrary coefficient cannot
enter this proposition just by being called a marked activity.

\subsection{The original small-class deletion preserves marked numerators}

On R-II p.378, the newly small selector implies that two specified old
characteristic factors are one. Denote them by $\Xi_{\rm old}$ and the new
factor by $\Xi_{\rm sm}$. Thus, on the original real configurations,
\begin{equation}
\Xi_{\rm sm}\Xi_{\rm old}=\Xi_{\rm sm}.
\label{eq:selector}
\end{equation}
Every other old restriction remains.

\begin{proposition}[Marked selector rewrite]
For an already defined original marked integrand $F$, the undivided
numerators with $\Xi_{\rm sm}\Xi_{\rm old}F$ and $\Xi_{\rm sm}F$ agree.
The equality holds coefficientwise for the complete finite-source polynomial,
including its zero-source coefficient. It therefore preserves the absolute
mass and does not enlarge the real set where $F$ is evaluated.
\end{proposition}
\begin{proof}
Multiply~\eqref{eq:selector} by $F$ before performing the original integrations.
For sharply constrained objects use the already defined joint distribution.
Equal full functionals have equal whole distributional derivatives. This
does not assert that an isolated derivative of $\Xi_{\rm old}$ is zero.
If product rules have exposed boundary terms, recombine all partners before
using the equality. No reference denominator is changed.
\end{proof}

\subsection{Gauge removal is an integrated equality}

The reverse gauge-fixing operation on the same page requires the prescribed
integration over $V_k$ on the newly small components. It is not equality of
fixed-$V_k$ conditional bodies. This qualification also applies to marks.

The directional force vector transforms as
$F_{e;\mu}(U^g)=\Ad_{g_x}F_{e;\mu}(U)$ for $e=(x,y)$.
Its invariant color pairings make the original $H$ source gauge invariant.
The original closed Wilson-word $C$ source has the same invariance. $V_H$ is
equivariant, and its complete divergence/contact packet, including
\eqref{eq:Casimir} and~\eqref{eq:b12}, respects that symmetry.
Carrying this invariant source polynomial through the original compact gauge
change of variables proves
\begin{equation}
\mathscr I_{\rm gf}[e^{A(\eta)}]
=\mathscr I_{\rm unfix}[e^{A(\eta)}]
\label{eq:gauge}
\end{equation}
after that required $V_k$ integration, with the original absolute
normalization on both sides. This proves equality of all complete source
coefficients and, when the common mass is nonzero, of their cumulants at
that integrated stage. It does not identify a conditional score, an individual
noninvariant current, or a fixed-field covariance. The boundary gauge needed
in the earlier joins cannot be removed early. Differentiating a parameterized
version of~\eqref{eq:gauge} differentiates its entire functional, including
every boundary and reference response.

\section{Connected estimates on the complete selected body}
\label{sec:connected}

Fix the actual common domain of the original complete operation and a positive
real reference body $\mu_0$. The history label can be one of its integration
variables. Do not replace a signed Mayer coefficient by a positive history law.
Use
\begin{equation}
T_zF=\int e^{\sigma_z(\xi)}F(z,\xi)\dd\mu_0(\xi),\quad
Z_0=\mu_0(1)>0,\quad Z_z=T_z1,\quad K_z=T_z/Z_z.
\label{eq:body}
\end{equation}
For the estimates in this section assume
\begin{equation}
\norm{\sigma_z}_\infty\le\eps<\log2,
\qquad b_\eps=2-e^\eps>0.
\label{eq:epsilon}
\end{equation}
The original small-coupling component comparison provides such a ceiling on
its specified class. It is not automatically uniform over an arbitrary
product of many components. Its explicit logarithmic and geometric factors
must be retained when using it quantitatively.

Since $|e^\sigma-1|\le e^\eps-1$,
\begin{equation}
b_\eps Z_0\le |Z_z|\le e^\eps Z_0.
\label{eq:nonzero}
\end{equation}
This lower bound is proportional to the genuine possibly tiny mass. No
order-one lower bound is assumed.

\begin{theorem}[Raw connected pair with one original reserve]
For $F,G\in L^2(\mu_0)$, with $F$ allowed to be Banach-valued,
\begin{equation}
\boxed{
\norm{\cC_z(F,G)}\le
\left(e^\eps+\frac{e^{2\eps}}{b_\eps}\right)
\bigl(\mu_0\norm{F}^2\,\mu_0|G|^2\bigr)^{1/2}.}
\label{eq:pairbound}
\end{equation}
If the original unused reserve obeys
\begin{equation}
\mu_0\bigl(e^{\theta E+\tau N}\bigr)\le B,
\qquad \theta,\tau>0,
\label{eq:reserve}
\end{equation}
and the actual envelopes are
\begin{equation}
\norm{F}\le A|\Gamma|W(1+E)^{D/2},\qquad
|G|\le g_kv(1+E)^{e/2},
\label{eq:envelopes}
\end{equation}
then
\begin{equation}
\boxed{\norm{\cC_z(F,G)}\le C_{\eps,\theta,D,e}
g_kvA|\Gamma|WB.}
\label{eq:onebank}
\end{equation}
Fixed carrier polynomials can be included using the $\tau N$ reserve.
\end{theorem}
\begin{proof}
The direct product term is bounded by $e^\eps$ times the product of the two
raw $L^2$ norms. For the quotient term, Cauchy--Schwarz gives
\[
\frac{|T_zF|\,|T_zG|}{|Z_z|}
\le\frac{e^{2\eps}}{b_\eps Z_0}
\bigl(\sqrt{Z_0}\sqrt{\mu_0\norm F^2}\bigr)
\bigl(\sqrt{Z_0}\sqrt{\mu_0|G|^2}\bigr).
\]
The mass cancels before any reserve is used. Bound each fixed energy power by
$C e^{\theta E}$ in the same original measure. The two square roots of its
one numerical majorant multiply to $B$, not $B^2$. This proves
\eqref{eq:onebank}. For a real positive scalar covariance, centered
Cauchy--Schwarz gives the sharper constant one in the unperturbed pair bound.
\end{proof}

The reserve in~\eqref{eq:reserve} is an input about the actual measure. This
theorem does not create it for a new history class. All engineering factors,
old coupling powers, and independently owned normalizers in $A,\Gamma,W,B$
remain visible.
These are bounds for the raw coefficient $Z\kappa$. A bound for the globally
normalized physical cumulant still needs the actual comparison between $B$
and that complete physical mass $Z$. Algebraic cancellation of intermediate
component inverses does not prove a cutoff-uniform estimate for $B/Z$.

\begin{theorem}[Every fixed finite source order]
For scalar marks $F_1,\ldots,F_r$ in $L^r(\mu_0)$,
\begin{align}
\left|Z_z\kappa_{K_z}(F_1,\ldots,F_r)\right|
&\le C_{r,\eps}\prod_{i=1}^r(\mu_0|F_i|^r)^{1/r},\notag\\
C_{r,\eps}
&=\sum_{\pi\vdash[r]}(|\pi|-1)!e^{|\pi|\eps}
b_\eps^{1-|\pi|}.
\label{eq:rbound}
\end{align}
If each mark has a fixed polynomial energy envelope
$A_i|\Gamma_i|W_i(1+E)^{d_i/2}$, one reserve yields
\begin{equation}
\left|Z_z\kappa_{K_z}(F_1,\ldots,F_r)\right|
\le C B\prod_i A_i|\Gamma_i|W_i.
\label{eq:rbank}
\end{equation}
\end{theorem}
\begin{proof}
For a partition $\pi$ with $p$ blocks, the raw cumulant term has denominator
$Z_z^{p-1}$. For each block $J$,
\[
\mu_0\prod_{i\in J}|F_i|
\le Z_0^{1-|J|/r}\prod_{i\in J}(\mu_0|F_i|^r)^{1/r}.
\]
The block product contains precisely $Z_0^{p-1}$, cancelling the denominator
power using~\eqref{eq:nonzero}. Summing the partition coefficients proves
\eqref{eq:rbound}. The $r$ moment roots of the common reserve give one $B$.
\end{proof}

For the nonlinear source action~\eqref{eq:finiteaction}, the exact formula is
\begin{equation}
\boxed{
[\eta_I]\log K_ze^{A(\eta)}
=\sum_{\pi\vdash I}\kappa_{K_z}(A_J:J\in\pi).}
\label{eq:nonlinearcumulants}
\end{equation}
Expand the cumulant series of $A$; the ordered block partitions cancel its
factorial. Consequently all inherited mixed coefficients are estimated with
their own saved grades and scalar factors. Replacing $A_J$ by products of
linear marks would change the transported source.

\section{Outer covariance and actual distinguished-interaction attachment}
\label{sec:outer}

\subsection{The complete joint measure includes covariance of means}

If an unchanged original positive outer factor genuinely integrates $y$, use
\begin{equation}
\dd\mu_0(y,\xi)=\dd\nu(y)\dd\mu_{0,y}(\xi).
\label{eq:joint}
\end{equation}
Under the uniform comparison and actual joint reserve, Section~\ref{sec:connected}
applies directly to this measure. Its normalization is the total mass, not
separately normalized inner fibers with unchanged outer weights.
Writing its disintegration as $\pi_0(\dd y)K_{0,y}(\dd\xi)$ gives
\begin{equation}
\Cov_{K_0}(F,G)
=\int\Cov_{K_{0,y}}(F,G)\dd\pi_0(y)
+\Cov_{\pi_0}(K_{0,y}F,K_{0,y}G).
\label{eq:totalcov}
\end{equation}
Both terms are in the joint estimate. If $y$ is still retained outside the
current operation, the estimate is only parameterwise until its original
outer integration is actually performed. A varying outer weight belongs to
the full score or to its boundary/reference distribution.

For any finite measure, the exact raw determinant identity is
\begin{equation}
Z\mu(FG)-\mu(F)\mu(G)
=\frac12\iint(F(x)-F(x'))(G(x)-G(x'))\dd\mu(x)\dd\mu(x').
\label{eq:twocopycov}
\end{equation}
For $\mu=\sum_\omega\mu_\omega$, its right side contains every ordered pair
$(\omega,\omega')$. Keeping only equal-history pairs loses the covariance
of conditional means.

\begin{corollary}[The protected common-row sector has a quantitative outer bound]
In the original PE12 sector, suppose the actual old mark factors through
the common retained rows and equals $R_j(y)$ before source localization.
Then
\begin{equation}
\cC_0(R_j,V_k)=Z_0\Cov_{\pi_0}(R_j,K_{0,y}V_k).
\label{eq:protectedouter}
\end{equation}
Whenever the already saved mark and the actual distinguished vertex satisfy
\eqref{eq:envelopes} on the same original joint measure,
\begin{equation}
\left\|Z_0\Cov_{\pi_0}(R_j,K_{0,y}V_k)\right\|
\le C g_kvA_j|\Gamma_j|W_{k,j}B.
\label{eq:protectedouterbound}
\end{equation}
\end{corollary}
\begin{proof}
Output measurability makes the inner covariance in~\eqref{eq:totalcov} zero.
It does not make the outer term zero. The raw pair estimate on the actual
joint measure proves~\eqref{eq:protectedouterbound}.
\end{proof}
This completes a numerical outer-covariance row for the verified factored
sector under its original joint envelope and reserve. It does not establish
those inputs for a touched reconstructed source. The common-row R step gives
no new regular age contraction; already earned weights keep their old labels.

\subsection{Outer complex comparison preserves the full covariance}

For fixed marks set $F^c=F-K_0F$, $G^c=G-K_0G$ and
\begin{equation}
w_z=\frac{e^{\sigma_z}}{K_0e^{\sigma_z}},\qquad
\delta_\eps=\frac{2(e^\eps-1)}{2-e^\eps}.
\end{equation}
Then $K_0w_z=1$ and $\norm{w_z-1}_\infty\le\delta_\eps$.
Direct expansion yields
\begin{align}
\Cov_{K_z}(F,G)-\Cov_{K_0}(F,G)
={}&K_0[(w_z-1)F^cG^c]\notag\\
&-K_0[(w_z-1)F^c]K_0[(w_z-1)G^c].
\label{eq:covdifference}
\end{align}
Consequently
\begin{equation}
Z_0|\Cov_{K_z}(F,G)-\Cov_{K_0}(F,G)|
\le(\delta_\eps+\delta_\eps^2)
(\mu_0|F^c|^2\mu_0|G^c|^2)^{1/2}.
\label{eq:covdifferencebound}
\end{equation}
The raw coefficient has an additional absolute-mass response:
\begin{align}
\cC_z-\cC_0
={}&Z_z(\Cov_{K_z}-\Cov_{K_0})
+(Z_z-Z_0)\Cov_{K_0},\notag\\
|\cC_z-\cC_0|
\le{}&\{e^\eps(\delta_\eps+\delta_\eps^2)+e^\eps-1\}
(\mu_0|F^c|^2\mu_0|G^c|^2)^{1/2}.
\label{eq:rawdifference}
\end{align}
The coefficient is $O(\eps)$ for $\eps\le1/4$. Moving marks require the
three additional covariance terms containing $\Delta F$, $\Delta G$, or
both, and the separate inherited $A_{12}$ response. The kernel comparison
does not supply those changed-readout bounds.

\subsection{A distinguished actual interaction gives the required coupling}

The outer analytic comparison $\sigma_z$ is not the interaction interpolation
used to prove MT36. For the latter use its actual bodies
$\dd\mu_t=e^{tV_k}\dd\mu_0$, when this is the specified interpolation.
For fixed marks,
\begin{align}
\partial_t\{Z_t\kappa_{K_t}(F_1,\ldots,F_r)\}
={}&Z_t\kappa_{K_t}(F_1,\ldots,F_r,V_k)\notag\\
&+(\mu_tV_k)\kappa_{K_t}(F_1,\ldots,F_r),
\label{eq:rawinterpolation}\\
\partial_t[\eta_I]\log K_te^{A(\eta)}
={}&\sum_{\pi\vdash I}\kappa_{K_t}(V_k,A_J:J\in\pi).
\label{eq:sourceinterpolation}
\end{align}
The second row of~\eqref{eq:rawinterpolation} is essential: the raw body
changes. Moving $A_J$ contribute their own differentiated terms.
If the actual real interpolation is positive and its moment reserve is
uniform in $t$, the same partition proof and
$|V_k|\le g_kv(1+E)^{e/2}$ give $O(g_k)$ raw responses with one reserve.
Complex source walls need their own verified comparison and admission.

The original clean inherited interaction has a useful established input.
The marginal-paired proper difference, localized with its original readout,
vanishes at $B=0$ for every allowed wall value (original IV10 and MT27--31).
On its admitted physical fluctuation disc, Schwarz gives
\begin{equation}
|V'_k(Y;U,J,B)|
\le C E_0g_k\norm B_{\infty,Y}e^{-\kappa_*d_k(Y)}.
\label{eq:actualvertex}
\end{equation}
This is the actual distinguished-vertex gain on that clean completed branch.
It does not put a normalized marked primitive in the unmarked $E_0$ class,
and it does not cover every mixed-R vertex automatically.

There are two legitimate ways to pay an explicit cutoff polynomial.
First, if the same actual comparison has a pointwise bound
$|\sigma_z|\le g_kvP(N,E)$, with the separate ceiling
$\norm\sigma\le1/4$, then
\begin{equation}
\partial_t\cC_t(F,G)
=\mu_t\bigl[\sigma_z(F-K_tF)(G-K_tG)\bigr],
\quad \dd\mu_t=e^{t\sigma_z}\dd\mu_0.
\label{eq:polynomialtilt}
\end{equation}
Expand the centered product and apply the same finite-partition H\"older
cancellation; the reserve pays $P$ and produces $C g_kB$ times the mark
amplitudes. A global supremum estimate alone does not prove this pointwise
polynomial hypothesis.

Second, if only $\eps_k\le Cg_kQ_k$ is known and a declared unused part
of the original mass bound is $Z_0\le B_ke^{-cP_k}$, then, for
$0<\delta<1$ with $\sup_kQ_ke^{-\delta cP_k}<\infty$,
\begin{equation}
\eps_kZ_0\le C_\delta g_kB_ke^{-(1-\delta)cP_k}.
\label{eq:bankpayment}
\end{equation}
The original R-II rare bound (1.89) supplies a possible bank on its own
class; its allocation must still be checked against payments already made.
All fixed polynomials are collected before spending this one fraction.
Neither method pays an unsupported old inverse-mesh power or creates lost
regular age.
Equation~\eqref{eq:bankpayment} directly pays a scalar or supremum envelope.
For energy-weighted marks, include a proved cutoff supremum in $Q_k$ or
use a corresponding moment-weighted unused bank. A bound on $Z_0$ alone
does not control $\mu_0(1+E)^D$.

\section{Saved age, grade, and the complete projection}
\label{sec:grade}

For a completed admitted stable mark born at $j$, put
$r=L^{j-k}$ and $\varrho=g_k/g_j$. The original test-functional norm is
$\cT=(C^{2,s})^*$ with its source anchor/path moments and lower/contact
filing. On the protected physical readout the MT44 envelopes take the form
\begin{equation}
\norm{F_{j\to k}}_{\cT}
\le C A_j|\Gamma_j|(w_2q^2+w_3q^3+w_4q^4),
\qquad q\le C(1+E)^{1/2}.
\label{eq:graded}
\end{equation}
For a clean fixed physical radius $\alpha$ the weights are
\begin{center}
\begin{tabular}{@{}c c c c@{}}
\toprule
Saved degree & $w_2$ & $w_3$ & $w_4$\\
\midrule
$d=2$ & $r^s\varrho^2$ & $r^\beta(g_k/\alpha)\varrho^2$
& $r^2(g_k/\alpha)^2\varrho^2$\\[3pt]
$d=3$ & $r^s\varrho^2$ & $r^\beta\varrho^3$
& $r^2(g_k/\alpha)\varrho^3$\\
\bottomrule
\end{tabular}
\end{center}
Here $\beta$ in $r^\beta$ is the original age exponent, not the Wilson action
coefficient. The table is used only after the full quadratic covariant
completion, its cubic path-area partner, and the prescribed lower rows have
been separated. R-origin marks use their actual shrinking radii and promotion
records, not a substituted clean $\alpha$.

\begin{corollary}[Connected attachment preserves each supplied saved row]
If the actual stored coefficient in the selected numerator satisfies
\eqref{eq:graded} on the same admitted joint domain as the distinguished
vertex and reserve, then
\begin{equation}
\boxed{
\norm{\cC_z(F_{j\to k},V_k)}_{\cT}
\le Cg_kvA_j|\Gamma_j|(w_2+w_3+w_4)B.}
\label{eq:gradedattachment}
\end{equation}
The same assertion holds for each already bounded current/contact row with
its own lowered grade.
\end{corollary}
\begin{proof}
Apply~\eqref{eq:onebank} to the three terms of the actual envelope, or test
against a unit $C^{2,s}$ function before applying the scalar proof. No age
factor is changed by this integration estimate.
\end{proof}

For $P_6$, Cauchy is taken on the actual common outer canonical disc. If the
physical radius is $\rho_k$ and canonical amplitude is $g_k$, an order-$m$
field derivative costs $(g_k/\rho_k)^m$. This ratio stays explicit until an
original radius estimate or an unused reserve controls it. The finite coarse
coordinate and anchor/path sums cost their actual carrier polynomial, paid
by the reserved locality exponent. On the verified fixed-radius completed
branch these are the original MT40--44 bounds. An unhealed live coefficient
does not receive a regular $P_6$ center.

Thus~\eqref{eq:gradedattachment} supplies the connected numerical step
required by MT36 once its incoming representative is admitted. It does not
prove that $R_j\circ\Phi_{j,k}$ in the actual mixed numerator satisfies
\eqref{eq:graded}. The distinction between this corollary and that missing
entry estimate is the central status distinction in this note.

\section{Full inherited-score and current response}
\label{sec:score}

\subsection{Outer derivatives of the actual complete body}

On a fixed-domain representation~\eqref{eq:body}, take commuting constant
outer-coordinate directions and set
\begin{equation}
S_a=\partial_a\sigma,\qquad L_a=S_a-K_zS_a.
\end{equation}
The score contains all varying inherited terms of that operator. Quotient
differentiation gives
\begin{align}
\partial_a(K_zF)&=K_zF_a+K_z(L_aF),\label{eq:firstscore}\\
\partial_{ab}(K_zF)&=K_zF_{ab}+K_z(L_aF_b)+K_z(L_bF_a)
+K_z(H_{ab}F),\label{eq:secondscore}\\
H_{ab}&=\sigma_{ab}-K_z\sigma_{ab}+L_aL_b
-\Cov_{K_z}(S_a,S_b).
\label{eq:Hscore}
\end{align}
In particular $K_zL_a=K_zH_{ab}=0$. Dropping the last normalizer response
in~\eqref{eq:Hscore} would already give a false formula for $F=1$.

For all finite orders, take the logarithm near one of $K_0e^\sigma$ and put
\begin{equation}
\ell=\sigma-\log K_0e^\sigma,\qquad
H_\alpha=e^{-\ell}\partial^\alpha e^\ell.
\end{equation}
Then
\begin{equation}
[\partial^\alpha,K_z]F
=\sum_{0<\delta\le\alpha}\binom{\alpha}{\delta}
K_z(H_\delta\partial^{\alpha-\delta}F),
\qquad K_zH_\delta=0\quad(\delta\ne0).
\label{eq:bell}
\end{equation}
For ordered outer vector fields with coefficients depending only on the
outer coordinates, use $H_{(X_m,\ldots,X_1)}=w^{-1}X_m\cdots X_1w$.
The recursion is $H_{(X,\mathbf X)}=XH_{\mathbf X}+(X\ell)H_{\mathbf X}$.
Product rules keep the original order and every coefficient derivative.
An actual output velocity depending on an integrated variable is not such
an outer field: it stays inside the marked numerator.

There is a fixed-direction Cauchy estimate with no ambient coordinate count.
Let $\norm\cdot_*$ be the actual Minkowski norm of an admitted outer ball
(including its different field and current radii). If $z$ has spare distance
$\eta>0$ in that norm, the polydisc
$z+\sum_{i=1}^m t_iu_i$, $|t_i|<\eta/(2m\norm{u_i}_*)$, is admitted.
Applying Cauchy to $w-1$, whose supremum is $O(\eps)$, gives
\begin{equation}
\norm{H_{u_m,\ldots,u_1}}_\infty
\le C_{\eps,m,\eta}\eps\prod_{i=1}^m\norm{u_i}_*.
\label{eq:directional}
\end{equation}
Multiplication by $w^{-1}$ is bounded by $e^{2\eps}$. Finite variable-field
words require their coefficient jets as well. Sums over bonds, colors,
anchors, or tensor indices retain their original Schur/tensor bounds.
Equation~\eqref{eq:directional} is not a bound for arbitrary internal Haar
derivatives of the inherited action.

\subsection{Boundary, reference, and absolute-mass responses}

If a parameter moves a selector, constraint, coarea factor, or reference,
write its actual complete tested derivative as
\begin{equation}
\partial_a(TF)=T(F_a)+T(S_aF)+\mathcal B_a(F).
\end{equation}
Then the genuine normalized response is
\begin{equation}
\partial_a(KF)=KF_a+\Cov_K(F,S_a)
+\frac{\mathcal B_a(F)-(KF)\mathcal B_a(1)}{Z}.
\label{eq:boundaryquotient}
\end{equation}
The full distribution is differentiated; exposed delta factors are not
multiplied separately. For commuting parameter derivatives the safe recursion is
\begin{equation}
\partial^\alpha(KF)=\frac1Z\left[
\partial^\alpha(TF)-\sum_{0<\delta\le\alpha}\binom{\alpha}{\delta}
(\partial^\delta Z)\partial^{\alpha-\delta}(KF)\right].
\label{eq:quotientrecursive}
\end{equation}
It retains every denominator partner.

For a smooth real variation $\dd\mu_t=\rho_t\dd\lambda$, let
$s_t=\partial_t\log\rho_t$, $F_t^c=F_t-K_tF_t$ and
$G_t^c=G_t-K_tG_t$. The raw variation is
\begin{equation}
\partial_t\cC_t(F_t,G_t)
=\mu_t(\dot F_tG_t^c+F_t^c\dot G_t+s_tF_t^cG_t^c).
\label{eq:rawscore}
\end{equation}
The score here is uncentered because the absolute body changes. A moving
sharp part adds $\langle\dot\mu_t^\partial,F_t^cG_t^c\rangle$.
The inherited contact adds the derivative of $\mu_t A_{12,t}$, including
the same boundary response. An independent outside factor adds
$a'\cC+a\partial_t\cC$. None of these terms follows from a native
Poincar\'e estimate.

\section{Connecting the verified one-step map to actual current numerators}
\label{sec:weakcurrent}

This section concerns precisely the original one-step averaging
map of~\cite[\S5]{PW} on its verified local plaquette sector. It does not identify a general
mixed R readout with that map. Its pivot inverse, local map jets, and two
Schur sums are bounded there by $C_{L,m}$. The source, pivot readsets,
frames, and original exterior data are the same as in~\cite[\S5]{PW}.

Let $\nu$ be the actual pre-pushforward selected joint measure or distribution,
including the inherited factor, coarea, selectors, and original constraints.
Variables not integrated in this operation remain parameters. In an admitted
output chart write the actual map as $q$ and $v_i^a=X_iq^a$.
The coefficient $v_i^a$ is generally a function of the integrated variables.
The weak estimates here use real field-variable tests and the real map. A
complex weight on this same real map can be bounded by absolute comparison.
A moving complex map requires a declared analytic test space; an arbitrary
$C^2$ test on real fields has no automatic holomorphic extension.


\begin{theorem}[Exact weak current pushforward]
Let
$L=\sum_{i,j}h_{ij}X_jX_i+\sum_i b_iX_i$, with the coefficients outside
the displayed derivatives, and define
$\langle L^*\nu,\psi\rangle=\langle\nu,L\psi\rangle$.
For a smooth admitted multiplier $F$,
\begin{align}
q_*(F L^*\nu)
={}&q_*((LF)\nu)-\partial_aq_*(C_F^a\nu)
+\partial_b\partial_aq_*(A_F^{ba}\nu),\label{eq:weakpush}\\
A_F^{ba}&=F\sum_{i,j}h_{ij}v_j^b v_i^a,\notag\\
C_F^a&=\sum_{i,j}h_{ij}\{(X_iF)v_j^a+(X_jF)v_i^a
+F X_jv_i^a\}+F\sum_i b_i v_i^a.
\label{eq:weakcoefficients}
\end{align}
Any additional original boundary/reference distribution $B$ contributes
$q_*(FB)$ without alteration.
\end{theorem}
\begin{proof}
Test against $\phi$ and expand $L(F\phi(q))$. The zeroth-order term is
$(LF)\phi(q)$. The first-order terms are exactly $C_F^a\partial_a\phi(q)$;
the second-order term is $A_F^{ba}\partial_b\partial_a\phi(q)$.
Move only these output test derivatives by the definition of distributional
derivative. This proves the signs in~\eqref{eq:weakpush}. If $L=V_2V_1$,
derivatives of the coefficients of $V_1$ are part of $b_i$. The complete
adjoint already contains any derivatives of $h_{ij},b_i$ required by its
original definition.
\end{proof}

The first-order version is
\begin{equation}
q_*(F D_V\nu)=q_*((VF)\nu)-\partial_aq_*(F(Vq^a)\nu).
\label{eq:firstweakpush}
\end{equation}
Thus the actual velocity stays inside the selected numerator. It is never
replaced by a fixed output vector before integration. These formulas hold
weakly with the complete sharp distribution, so no score of a disconnected
selected body and no native convexity assumption are needed.

\subsection{A concrete first/second $H$ current estimate}

Use the original physical vectors $V_i=V_H(f_i)$ on one source lattice with
the same action coefficient $\beta_s\ge1$. Let $E_Y=\beta_s\sum_{p\in Y}s_p$
for the literal union of source, compensating-pivot, and readout stars. The
map and Wilson Hessian have bounded local incidence on this verified sector.
Compact Wilson force bounds and the Schur estimates give
\begin{equation}
\norm{dq\,V_i}_2
\le C_L|\sigma_i|\beta_s^{-3/2}\norm{f_i}_\infty E_Y^{1/2}.
\label{eq:Hvelocity}
\end{equation}
For $F=1$ and $L=V_2V_1$, the second tensor and first drift are
$A=(dqV_2)\otimes(dqV_1)$ and $C=V_2(V_1q)$. Their bounds are
\begin{align}
\norm A_{\rm proj}
&\le C_L|\sigma_1\sigma_2|\beta_s^{-3}
\norm{f_1}_\infty\norm{f_2}_\infty E_Y,\label{eq:Htensor}\\
\norm C_2
&\le C_L|\sigma_1\sigma_2|\norm{f_1}_\infty\norm{f_2}_\infty
\{\beta_s^{-3}E_Y+\beta_s^{-5/2}E_Y^{1/2}\}.
\label{eq:Hdrift}
\end{align}
The projective tensor norm pairs with the operator norm of the output
Hessian. The second term in~\eqref{eq:Hdrift} is the Wilson force derivative,
not a discarded lower row.

To check these bounds, the force vector has squared norm at most
$C\beta_s^{-1}E_Y$. The Wilson Hessian is a sparse bounded matrix, hence
has a uniform $\ell^2$ operator bound. The finite-range second derivative of
the original map has a bounded $\ell^2\times\ell^2\to\ell^2$ norm on
this chart. Applying these statements to $V_2(V_1q)$ gives respectively the
$\beta_s^{-5/2}\sqrt{E_Y}$ and $\beta_s^{-3}E_Y$ terms. No claim about an
arbitrary inherited-prefix derivative is used.

In admitted canonical output coordinates of amplitude $h$, the tensor and
drift acquire $h^{-2}$ and $h^{-1}$, respectively; the actual logarithm-chart
jets are retained. To display the saved normalization explicitly, write
$\beta_s=c_\beta\gamma_s^{-2}$, $\sigma_i=\gamma_i^{-2}a^{-2}$,
$R_i=(\gamma_s/\gamma_i)^2$, and $t=\gamma_s/h$. Up to the fixed powers
of $c_\beta$, the first-current prefactor is $a^{-2}R_it$. The second
tensor, drift-energy, and drift-square-root prefactors are, respectively,
\begin{equation}
a^{-4}R_1R_2t^2,\qquad
a^{-4}R_1R_2\gamma_st,\qquad
a^{-4}R_1R_2t.
\label{eq:savedcurrentratios}
\end{equation}
No $\gamma_i$ is reset to $h$. These formulas apply to the displayed
Wilson-force vectors; a general inherited vector must first have its actual
representation and coefficient bounds established. Only in the explicitly
matched specialization
\begin{equation}
h=\gamma,\qquad \beta_s=c_\beta\gamma^{-2},\qquad
\sigma_i=\gamma^{-2}a^{-2},\qquad 0<\gamma\le1
\label{eq:matchedcurrent}
\end{equation}
do both rows have the bound
\begin{equation}
\norm{A^{\rm can}}_{\rm proj}+\norm{C^{\rm can}}_2
\le C_{L,c_\beta}a^{-4}\norm{f_1}_\infty\norm{f_2}_\infty
(E_Y+E_Y^{1/2}).
\label{eq:physicalcurrentbound}
\end{equation}
For a positive actual joint measure satisfying
$\nu(e^{\theta E_Y})\le B$, use the field-test norm
\begin{equation}
\norm\phi_{C^2_{\rm out}}
=\max\{\norm\phi_\infty,\sup\norm{\nabla\phi}_2,
\sup\norm{D^2\phi}_{2\to2}\}.
\label{eq:outputtestnorm}
\end{equation}
Testing against a unit function in this norm gives
\begin{equation}
\norm{q_*( (V_2V_1)^*\nu)}_{(C^2_{\rm out})^*}
\le C_{L,c_\beta,\theta}a^{-4}
\norm{f_1}_\infty\norm{f_2}_\infty B.
\label{eq:weakcurrentbound}
\end{equation}
The factor $a^{-4}$ remains. The bank contains the literal energy and carrier
count. This weak field-variable norm is distinct from the source-test norm
$\cT=(C^{2,s})^*$ in Section~\ref{sec:grade}.

For inherited $F$, every $LF$ and $X_iF$ term in~\eqref{eq:weakcoefficients}
must have its actual admitted grade. To estimate connected coefficients,
first form the actual joint finite-source coefficient, with every ordered
lower/contact term. Then test its complete weak representation and use the
resulting ordinary joint coefficient functions in Section~\ref{sec:connected}.
Separately testing individual currents and multiplying their integrands does
not reconstruct the original contact polynomial. Do not multiply
distributional delta rows.
This is a genuine restricted current-transfer estimate for the actual
selected body; it does not bound the full original $H$ source by itself.

\subsection{The inherited-score correction is still part of $H$}

For an inherited smooth factor $G$ and the original joint sharp factor $\Xi$,
the exact identity is
\begin{equation}
\OH\,\Xi G\rho_s
=r_H\Xi G\rho_s+D_{V_H}(\Xi G\rho_s)
+\rho_s V_H(\Xi G).
\label{eq:historycorrection}
\end{equation}
The last term includes $+V_H\log G$ in its smooth ordinary part and all
selector/constraint responses. The weak estimate controls the complete
current term in~\eqref{eq:historycorrection}; it does not set the other terms
to zero or bound the unknown internal derivative of $G$. Coarea and moving
reference responses enter the same full product rule after a coordinate
change. In particular, the $a^{-4}$ physical second-current row, the Casimir
action row, and the complete $b_{12}$ must be transported together.

\section{The touched stable readout: exact endpoint and unproved domain}
\label{sec:missing}

The actual first-entry physical representative is LM42, evaluated with its
own current. For the fixed original determining problem the candidate
source-only wall extension, wherever admitted, is
\begin{equation}
\widehat\Gamma_{\omega,s}=P_{\rm old}q_{\rm old}
\Gamma^R_{\omega,s},\qquad
\widehat J_{\omega,s}=J_{\rm phys}(\widehat\Gamma_{\omega,s}).
\label{eq:actualreadout}
\end{equation}
The same global determining problem, exterior data, and complete packet
collar are used for every root. At the genuine admitted physical endpoint,
uniqueness gives the original readout; LM46a also identifies all physical
tangent derivatives and complete weak contact identities there. The physical
numerator is therefore not alleged to be off its own-current graph.

The missing assertion is uniform control of~\eqref{eq:actualreadout} on the
old real selector support and every simultaneous preparation/source wall:
\begin{equation}
(\widehat\Gamma_{\omega,s},\widehat J_{\omega,s})
\in\mathcal C^{\rm spare}_{j,k},
\quad\text{with its required finite jets and anchor sums.}
\label{eq:admissiontarget}
\end{equation}
The domain on the right is the one on which the stored representative's
graded bound was proved. Native member admission in R-II (1.65) is not this
statement. The invertible local pivot of~\cite[\S5]{PW} also does not bound the old target
$q_{\rm old}\Gamma^R$ in the requisite norm or supply fractional-curvature
and actual-current bounds across a rough join.

The exact numerical failure to be repaired is now explicit. If only a
native bound $b_j$ with the full birth-core cost is available on the actual
record, the raw connected estimate gives
$C g_kb_jB$, with its full old coupling cost. The required row is
$C g_kA_j|\Gamma_j|(w_2+w_3+w_4)B$. Passing between these two statements
requires~\eqref{eq:admissiontarget} and the actual readout identification;
neither denominator cancellation nor an eventual rare factor supplies it.

\subsection{A quantitative defect formula, when its segment is admitted}

On one common paired chart let
$\iota(U)=(U,J_{\rm phys}(U))$ and
$\Delta(U,J)=J-J_{\rm phys}(U)$. For the same stored analytic $F$, assume
the entire segment $\gamma_t=(U,J_{\rm phys}(U)+t\Delta)$ and a transverse
complex margin $r_J$ are admitted. Then
\begin{align}
F(U,J)-F(U,J_{\rm phys}(U))
&=\int_0^1D_JF(\gamma_t)[\Delta]\dd t,\label{eq:Hadamard}\\
|F(U,J)-F(U,J_{\rm phys}(U))|
&\le C\norm F_{\rm tube}\frac{\norm\Delta}{r_J}.
\label{eq:Hadamardbound}
\end{align}
For an actual readout parameter $y$, its derivative is
\begin{equation}
D_y\!\int_0^1D_JF(\gamma_t)[\Delta]\dd t
=\int_0^1\{D(D_JF)(\gamma_t)[D_y\gamma_t,\Delta]
+D_JF(\gamma_t)[D_y\Delta]\}\dd t.
\label{eq:defectjet}
\end{equation}
Higher terms contain a derivative of $\Delta$, possibly its zeroth
derivative. Thus a successful defect route must control the actual mismatch
jets, readout jets, and common tube with sufficient saved weights. If
$\norm F_{\rm tube}$ is only the large native norm, that factor persists.
These formulas give a precise possible estimate; no claim is made that the
actual mixed readout has satisfied their hypotheses.

\subsection{What exact physical-ideal regrouping can and cannot do}

The original XP34--35a already separates the physical graph ideal
$\mathcal I_{\rm phys}=\{F:\iota^*F=0\}$ from independent-current values.
A graph-vanishing coefficient need not vanish at independently weakened
walls. For example, for $J=u^2$, the function $J-u^2$ composed with
$(u,J)\mapsto(u,sJ)$ becomes $(s-1)u^2$ on the physical graph.
This is a diagnostic of an inference, not a counterexample satisfying the
actual R selectors.

Let $\mathscr S_k$ synthesize the complete localized packet into its physical
numerator. If exact source-free transport obeys
$\mathscr S_{k+1}\mathscr R=\mathcal K_k^{\rm phys}\mathscr S_k$, then
$\mathscr R(\ker\mathscr S_k)\subseteq\ker\mathscr S_{k+1}$.
This algebraic cancellation is valid before estimates. For a useful local
connected cancellation one needs pointwise zero before multiplying other
integrated marks, or coefficientwise zero of the entire relevant finite-source
packet on every retained record. Merely $K_yF=0$ does not imply
$K_y(FG)=0$ for another internal mark $G$. Such conditional zero survives
multiplication only by retained-measurable partners without further input.

Even exact complete-packet cancellation does not prove a bounded local
section of physical restriction in the original carrier norm. A simple
analytic check isolates the issue: on $|u|,|J|<1$ with graph $J=u^2$, take
$f_n(u)=(u/\rho)^n$, $0<\rho<1$. Its norm on $|u|<\rho$ is one and it
vanishes to any fixed finite jet order for large $n$. An analytic paired
extension agreeing there must agree on the entire graph by the identity
theorem. At any $\rho<s<1$ its paired norm is at least $(s/\rho)^n$.
Quotienting by the graph ideal cannot remove this domain loss. This example
does not disprove actual admission; it rules out deriving it from a smaller
physical norm and an algebraic quotient alone.

\section{Testing the conditional contact prescription in the nonlinear law}
\label{sec:contact}

The explicit prescription of~\cite[\S3]{PW} for a sparse independent family
$\mathcal F$ subtracts a bath-dependent polynomial $Q_{\mathcal F}$ of
source degree at least two and satisfies
\begin{equation}
K_{\mathcal F}\exp\!\left(\sum_i\eta_i\OH(f_i)
-Q_{\mathcal F}(\eta)\right)
=\exp\!\left(\sum_i\eta_i m_i\right),
\qquad m_i=K_{\mathcal F}\OH(f_i).
\label{eq:conditionalprescription}
\end{equation}
This changes the source hierarchy and leaves its projected covariance.
It preserves the zero-source Wilson mass. The original source has the
additional factor $e^{Q_{\mathcal F}}$ and is kept distinct.

\begin{theorem}[A legal sparse prescription still leaves a non-tight projection]
Take a real test $f$ equal to one on a fixed positive-volume box. Choose the
sparse witnessing links of~\cite[\S2.1]{ND} with mutually disjoint complete H determining,
plaquette-star, and good-bath patches, enlarging their fixed separation if
needed. Split them into two positive-density subfamilies
$\mathcal F,\mathcal G$ inside the plateau. Eliminate only $\mathcal F$.
Use the original birth normalization $\beta_s=\beta=c_\beta\gamma^{-2}$
and $\beta\ge\beta_0$. For the exact nonlinear Wilson law on compatible
sufficiently fine even tori,
\begin{align}
\Var(K_{\mathcal F}\OH(f))&\ge c_fa^{-8},\label{eq:residualvariance}\\
\sup_c\Pr\{|K_{\mathcal F}\OH(f)-c|\le L\}
&\le\min\{1,C_f(e^{-c'_f\beta}+a^2+a^4L)\}.
\label{eq:residualdispersion}
\end{align}
The constants have the same fixed plateau/patch-fit scope as in~\cite[\S2]{ND} and~\cite[\S2]{PW}.
\end{theorem}
\begin{proof}
Assign every original H summand reading $U_g$, $g\in\mathcal G$, to its
complete $h_g(U_g;b)$, and write
$\OH=O_0+\sum_{g\in\mathcal G}h_g$. Stencil separation makes $O_0$
independent of every $U_g$ and each $h_g$ independent of the links in
$\mathcal F$. No plaquette containing a link of $\mathcal F$ reads a
link of $\mathcal G$. Hence the exact conditional kernel $K_{\mathcal F}$
and its one-link denominators are independent of the $\mathcal G$ links,
and
\begin{equation}
K_{\mathcal F}\OH=K_{\mathcal F}O_0
+\sum_{g\in\mathcal G}h_g(U_g;b).
\label{eq:unchangedG}
\end{equation}
Condition on the remaining bath after marginalizing $\mathcal F$. The
$\mathcal G$ variables still have their original independent one-link
Wilson laws: all integrated $\mathcal F$ factors are bath-only for them.
Their good-bath probabilities are the original ones under marginalization.
The complete one-link variance lower bound of~\cite[\S2.3]{ND} applies to every good $g$.
There are $c_fa^{-4}$ such possible links, and the physical normalization
contributes $a^{-4}$ to each variance after its $\beta$ factors cancel.
Taking the conditional variance proves~\eqref{eq:residualvariance}; no
independence of good-bath events is needed. Applying the same conditional
concentration argument as in~\cite[\S2.4]{ND}, with the improved exponential
bad-bath bound of~\cite[\S2]{PW}, to the positive-density family $\mathcal G$ proves
\eqref{eq:residualdispersion}.
\end{proof}

Along $\beta(a)\to\infty$, the right side of~\eqref{eq:residualdispersion}
tends to zero for fixed $L$, even after arbitrary deterministic centering.
For a sufficiently large fixed $\beta$, its limiting bound is less than a
fixed number below one, which already rules out tightness. Variance divergence
alone would not justify this conclusion; the concentration estimate does.

This proves failure of this specified legal one-pass sparse choice. It is
not a no-go theorem for every possible multiscale contact prescription.
In particular, a further subtraction must be derived and tested on the
remaining projected hierarchy; the exact identity
\eqref{eq:conditionalprescription} is not by itself a renormalization proof.

\section{S94 and the remaining quantitative targets}
\label{sec:status}

No new ultraviolet-uniform pressure theorem is claimed. The exact
Wilson comparison of~\cite[\S8]{PW} retains a cost proportional to
$a^{-3}[C_G+\frac12\dim(G)\log(1+\beta)]/t_*$, which is insufficient.
Its precise two-copy target remains
\begin{equation}
\boxed{
\int K(U,V)K(V,U)\dd U\dd V
\ge e^{-C_{t_*}\Vol(\Sigma)}
\left(\int K(U,U)\dd U\right)^2,}
\label{eq:S94}
\end{equation}
for the original physical slab kernel
$K=(TP)^{t_*/a}$ on the product-Haar boundary space. The constant must be
independent of the ultraviolet spacing in the stated physical regime.
Both sides use the same original normalization. If the kernel is written
as an actual nonnegative history sum, the left side contains every pair of
histories, including unequal ones. Section~\ref{sec:connected} does not
bound this global two-copy overlap or the required excitation trace.

\begin{center}
\small
\begin{tabularx}{\linewidth}{@{}>{\raggedright\arraybackslash}p{0.25\linewidth}X@{}}
\toprule
Interface & Status\\
\midrule
Certified native marked evaluation & Carried on its original domain with its
full birth-core cost; no conversion to a stable age norm.\\[3pt]
Small-class marked rewrite & Exact selected numerator identity and complete
distributional response; only the proved redundant factors are removed.\\[3pt]
Gauge removal & Full source equality after the prescribed integration; no
fixed-retained-field body or score equality.\\[3pt]
Complete-body connected estimates & Proved at every fixed source order, with
one actual reserve and the genuine mass cancellation.\\[3pt]
Protected outer covariance & Quantitative $g_k$ attachment bound on the actual
joint measure when its already verified envelope and reserve apply.\\[3pt]
Outer score/current response & Exact Bell and quotient identities; quantitative
fixed-direction bounds on the admitted outer domain.\\[3pt]
One-step physical weak currents & Exact selected-measure pushforward and
first/second current bounds on the verified map of~\cite[\S5]{PW}; physical $a^{-4}$ retained.\\[3pt]
Touched stable admission & Open: actual old-target reconstruction, simultaneous
walls, spare physical/current domain, finite jets and source anchor sums.\\[3pt]
All-history MT36 & Open: the missing input above, inherited ordinary score
corrections, full current/contact grades, and actual distinguished-vertex inputs
must hold together.\\[3pt]
Conditional contact prescription & Exact changed hierarchy; a specified legal
one-pass choice has a proved non-tight nonlinear projection.\\[3pt]
S94 & The ultraviolet-uniform two-copy slab inequality is still unproved.\\
\bottomrule
\end{tabularx}
\end{center}

The next main estimate is not another native spectral-gap statement. It is
\eqref{eq:admissiontarget} for the literal touched old readout, with a strict
margin on the simultaneous mixed wall domain and its complete current jets.
The exact R-I ranges and selectors would permit a literal check of the
native geometry, but they would not themselves prove the stronger
marked estimate. On a successful entry, the raw connected calculus above
already gives the required numerical attachment without discarding the
outer covariance or dividing by a small selected-history mass.

\appendix
\section{Source map}

The arguments above retain (i) the exact
$Z_0^{p-1}$ H\"older cancellation in each cumulant partition,
(ii) the absolute-body term in raw interpolation and covariance derivatives,
(iii) the signs and lower rows in the weak second-current pushforward,
(iv) the integrated scope of gauge equivalence, and (v) the unchanged
one-link conditional law in the nonlinear residual-contact test.
Two shortcuts are not available:
an integrated-variable output velocity cannot be treated as an outer vector
field, and a one-mark conditional zero cannot erase a product with another
integrated mark. Native Poincar\'e domains are not enlarged.

\begin{center}
\small
\begin{tabularx}{\linewidth}{@{}>{\raggedright\arraybackslash}p{0.35\linewidth}X@{}}
\toprule
Location & Use here\\
\midrule
R-II (1.65), pp.373--375 & Original native prepared-domain inclusion after
three determining-set joins.\\[3pt]
R-II p.378; (1.71)--(1.75), pp.379--380 & Small-class rewrite, integrated gauge
equivalence, and actual complete-body comparison.\\[3pt]
R-II (1.89), p.387 & Original component rare factor; an allocation must still
be reserved before using it for a new cost.\\[3pt]
\cite[pp.~686--688]{W} MT21--25a & Primitive repaired readout, separate canonical
and full birth-core norms.\\[3pt]
\cite[pp.~678--680, 689--696]{W} IV1--12; MT27--48 & Actual clean distinguished-vertex
gain, MT36 target, graded table, proper-difference disc.\\[3pt]
\cite[pp.~726--727]{W} XP34--35a & Physical restriction ideal and its wall/naturality
limitations.\\[3pt]
\cite[pp.~778--783]{W} LM31--48 & Completed regular R birth domain, actual-current
branch, physical endpoint, and unresolved protected entry.\\[3pt]
\cite[pp.~932--935]{W} AE8.1--8.10 & Literal prepared composite versus old physical
readout; necessary assigned-scale curvature test.\\[3pt]
\cite[pp.~1174--1180]{W} PE12 & Common-row factorization and its limited inner
covariance cancellation.\\[3pt]
\cite{PW}, \S\S1--3, 7--9, 10--11 & Casimir/contact rows, improved Wilson
moments, one-step map, native comparison scope, and S94 target.\\
\bottomrule
\end{tabularx}
\end{center}

\begingroup\small\raggedright
\begin{thebibliography}{9}
\setlength{\itemsep}{3pt}
\bibitem{W} Wilson-source estimates, volumes 1--10 (\texttt{lattice\_theory/source\_estimates} in this collection), cited by source page.

\bibitem{ND} \emph{Nonlinear dispersion and unequal spatial weights}, \texttt{lattice\_theory/block\_and\_transfer/conditional\_estimates/nonlinear\_dispersion\_and\_unequal\_spatial\_weights.tex} in this collection. Nonlinear bare-$H$ dispersion and
the restricted native estimates.
\bibitem{PW} \emph{Physical Wilson contacts and local conditional subtraction}, \texttt{lattice\_theory/block\_and\_transfer/conditional\_estimates/physical\_wilson\_contacts\_and\_local\_conditional\_subtraction.tex} in this collection.
Local conditional contact prescription, actual one-step current lifts,
same-record native comparison, and the two-copy pressure target.
\bibitem{balabanRII} T. Balaban, \emph{Large field renormalization. II.
Localization, exponentiation, and bounds for the R operation},
Communications in Mathematical Physics \textbf{122} (1989), 355--392.
\href{https://doi.org/10.1007/BF01238433}{doi:10.1007/BF01238433}.
\bibitem{balabanRI} T. Balaban, \emph{Large field renormalization. I.
The basic step of the R operation}, Communications in Mathematical Physics
\textbf{122} (1989), 175--202.
\href{https://doi.org/10.1007/BF01257412}{doi:10.1007/BF01257412}.
Cited for identification only; its defining ranges are not used here.
\end{thebibliography}
\endgroup
\end{document}
